% GENERATED FILE. Edit canonical lessons, then run npm run generate:books.
% canonical-source-hash: fnv1a64:17318ce6355140ab
% canonical-lessons: PA-C356-kaimra, PA-C356-tasvir, PA-C356-foto, PA-C356-ailbam, PA-C356-farem

\chapter{Pictures and Photos}
\label{ch:pa-a2-pictures-and-photos}
\hlchaptermodality{356}

\begin{chapteropening}
\textbf{By the end of this chapter:} \emph{I can say which pictures and photos I like.}

Complete the last lesson of chapter 356: I can say which pictures and photos I like.
\end{chapteropening}

\section[\texorpdfstring{\pa{ਕੈਮਰਾ} (kaimrā)}{ਕੈਮਰਾ (kaimrā)}]{\pa{ਕੈਮਰਾ} (kaimrā) — a camera}

\label{lesson:PA-C356-kaimra}



\begin{quote}
Before the new one: say the Punjabi for to make someone laugh, then the Punjabi for to tease.
\end{quote}

\subsection*{You'll want to know: \pa{ਕੈਮਰਾ}}
\textbf{\pa{ਕੈਮਰਾ}} — \emph{kaimrā} — \textquotedblleft{}a camera\textquotedblright{}.

\begin{quote}
\pa{ਮੇਰੇ} \pa{ਕੋਲ} \pa{ਇੱਕ} \pa{ਪੁਰਾਣਾ} \pa{ਕੈਮਰਾ} \pa{ਹੈ।}
\end{quote}

\begin{quote}
\emph{mere kol ikk purāṇā kaimrā hai.}
\end{quote}

\begin{quote}
I have an old camera.
\end{quote}

A masculine noun.

\begin{cousinweb}
\textbf{\pa{ਕੈਮਰਾ}} (\emph{kaimrā}) is borrowed from English \emph{camera}.
\end{cousinweb}

\begin{grammarlens}[title={What you've built}]
One more word for this chapter.
\end{grammarlens}

\subsection*{Your turn}
\begin{itemize}
\raggedright
  \item \emph{Say it:} \emph{kaimrā}
  \item \emph{Say it:} \emph{kaimrā}, once more
  \item \emph{Say it:} \emph{mere kol ikk purāṇā kaimrā hai.} — I have an old camera.
  \item \emph{Recall:} say the Punjabi for to make someone laugh, then the Punjabi for to tease, then say \emph{kaimrā} again
\end{itemize}

\begin{tcolorbox}[breakable,colback=teal!4,colframe=teal!35!black,title={Before you move on}]
What does \pa{ਕੈਮਰਾ} mean? (\textquotedblleft{}A camera\textquotedblright{}.) Say it once more.
\end{tcolorbox}
\section[\texorpdfstring{\pa{ਤਸਵੀਰ} (tasvīr)}{ਤਸਵੀਰ (tasvīr)}]{\pa{ਤਸਵੀਰ} (tasvīr) — a picture}

\label{lesson:PA-C356-tasvir}



\begin{quote}
Before the new one: say the Punjabi for to tease, then the Punjabi for a camera.
\end{quote}

\subsection*{You'll want to know: \pa{ਤਸਵੀਰ}}
\textbf{\pa{ਤਸਵੀਰ}} — \emph{tasvīr} — \textquotedblleft{}a picture\textquotedblright{}.

\begin{quote}
\pa{ਮੈਨੂੰ} \pa{ਇਹ} \pa{ਤਸਵੀਰ} \pa{ਬਹੁਤ} \pa{ਪਸੰਦ} \pa{ਹੈ।}
\end{quote}

\begin{quote}
\emph{mainū̃ eh tasvīr bahut pasand hai.}
\end{quote}

\begin{quote}
I like this picture very much.
\end{quote}

A feminine noun.

\begin{cousinweb}
\textbf{\pa{ਤਸਵੀਰ}} (\emph{tasvīr}) is a loan, usually traced through Persian to Arabic \emph{taṣwīr}, \textquotedblleft{}picturing\textquotedblright{}.
\end{cousinweb}

\begin{grammarlens}[title={What you've built}]
One more word for this chapter.
\end{grammarlens}

\subsection*{Your turn}
\begin{itemize}
\raggedright
  \item \emph{Say it:} \emph{tasvīr}
  \item \emph{Say it:} \emph{tasvīr}, once more
  \item \emph{Say it:} \emph{mainū̃ eh tasvīr bahut pasand hai.} — I like this picture very much.
  \item \emph{Recall:} say the Punjabi for to tease, then the Punjabi for a camera, then say \emph{tasvīr} again
  \item \emph{Write it:} copy \textbf{\pa{ਤਸਵੀਰ}} once, checking the whole word
\end{itemize}

\begin{tcolorbox}[breakable,colback=teal!4,colframe=teal!35!black,title={Before you move on}]
What does \pa{ਤਸਵੀਰ} mean? (\textquotedblleft{}A picture\textquotedblright{}.) Say it once more.
\end{tcolorbox}
\section[\texorpdfstring{\pa{ਫ਼ੋਟੋ} (foṭo)}{ਫ਼ੋਟੋ (foṭo)}]{\pa{ਫ਼ੋਟੋ} (foṭo) — a photograph}

\label{lesson:PA-C356-foto}



\begin{quote}
Before the new one: say the Punjabi for a camera, then the Punjabi for a picture.
\end{quote}

\subsection*{You'll want to know: \pa{ਫ਼ੋਟੋ}}
\textbf{\pa{ਫ਼ੋਟੋ}} — \emph{foṭo} — \textquotedblleft{}a photograph\textquotedblright{}.

\begin{quote}
\pa{ਮੇਰਾ} \pa{ਭਰਾ} \pa{ਕੈਮਰੇ} \pa{ਨਾਲ} \pa{ਪਰਿਵਾਰ} \pa{ਦੀ} \pa{ਫ਼ੋਟੋ} \pa{ਖਿੱਚਦਾ} \pa{ਹੈ।}
\end{quote}

\begin{quote}
\emph{merā bharā kaimre nāl parivār dī foṭo khichchdā hai.}
\end{quote}

\begin{quote}
My brother takes a photo of the family with the camera.
\end{quote}

A feminine noun.

\begin{cousinweb}
\textbf{\pa{ਫ਼ੋਟੋ}} (\emph{foṭo}) is borrowed from English \emph{photo}.
\end{cousinweb}

\begin{grammarlens}[title={What you've built}]
One more word for this chapter.
\end{grammarlens}

\subsection*{Your turn}
\begin{itemize}
\raggedright
  \item \emph{Say it:} \emph{foṭo}
  \item \emph{Say it:} \emph{foṭo}, once more
  \item \emph{Say it:} \emph{merā bharā kaimre nāl parivār dī foṭo khichchdā hai.} — My brother takes a photo of the family with the camera.
  \item \emph{Recall:} say the Punjabi for a camera, then the Punjabi for a picture, then say \emph{foṭo} again
  \item \emph{Write it:} copy \textbf{\pa{ਫ਼ੋਟੋ}} once, checking the whole word
\end{itemize}

\begin{tcolorbox}[breakable,colback=teal!4,colframe=teal!35!black,title={Before you move on}]
What does \pa{ਫ਼ੋਟੋ} mean? (\textquotedblleft{}A photograph\textquotedblright{}.) Say it once more.
\end{tcolorbox}
\section[\texorpdfstring{\pa{ਐਲਬਮ} (ailbam)}{ਐਲਬਮ (ailbam)}]{\pa{ਐਲਬਮ} (ailbam) — a photo album}

\label{lesson:PA-C356-ailbam}



\begin{quote}
Before the new one: say the Punjabi for a picture, then the Punjabi for a photograph.
\end{quote}

\subsection*{You'll want to know: \pa{ਐਲਬਮ}}
\textbf{\pa{ਐਲਬਮ}} — \emph{ailbam} — \textquotedblleft{}a photo album\textquotedblright{}.

\begin{quote}
\pa{ਸਾਡੀਆਂ} \pa{ਫ਼ੋਟੋਆਂ} \pa{ਐਲਬਮ} \pa{ਵਿੱਚ} \pa{ਹਨ।}
\end{quote}

\begin{quote}
\emph{sāḍīāṁ foṭoāṁ ailbam vich han.}
\end{quote}

\begin{quote}
Our photos are in the album.
\end{quote}

A masculine noun.

\begin{cousinweb}
\textbf{\pa{ਐਲਬਮ}} (\emph{ailbam}) is borrowed from English \emph{album}.
\end{cousinweb}

\begin{grammarlens}[title={What you've built}]
One more word for this chapter.
\end{grammarlens}

\subsection*{Your turn}
\begin{itemize}
\raggedright
  \item \emph{Say it:} \emph{ailbam}
  \item \emph{Say it:} \emph{ailbam}, once more
  \item \emph{Say it:} \emph{sāḍīāṁ foṭoāṁ ailbam vich han.} — Our photos are in the album.
  \item \emph{Recall:} say the Punjabi for a picture, then the Punjabi for a photograph, then say \emph{ailbam} again
  \item \emph{Write it:} copy \textbf{\pa{ਐਲਬਮ}} once, checking the whole word
\end{itemize}

\begin{tcolorbox}[breakable,colback=teal!4,colframe=teal!35!black,title={Before you move on}]
What does \pa{ਐਲਬਮ} mean? (\textquotedblleft{}A photo album\textquotedblright{}.) Say it once more.
\end{tcolorbox}
\section[\texorpdfstring{\pa{ਫ਼ਰੇਮ} (farem)}{ਫ਼ਰੇਮ (farem)}]{\pa{ਫ਼ਰੇਮ} (farem) — a picture frame}

\label{lesson:PA-C356-farem}



\begin{quote}
Before the new one: say the Punjabi for a photograph, then the Punjabi for a photo album.
\end{quote}

\subsection*{You'll want to know: \pa{ਫ਼ਰੇਮ}}
\textbf{\pa{ਫ਼ਰੇਮ}} — \emph{farem} — \textquotedblleft{}a picture frame\textquotedblright{}.

\begin{quote}
\pa{ਤਸਵੀਰ} \pa{ਲਈ} \pa{ਇੱਕ} \pa{ਨਵਾਂ} \pa{ਫ਼ਰੇਮ} \pa{ਖ਼ਰੀਦੋ।}
\end{quote}

\begin{quote}
\emph{tasvīr laī ikk navāṁ farem kharīdo.}
\end{quote}

\begin{quote}
Buy a new frame for the picture.
\end{quote}

A masculine noun.

\begin{cousinweb}
\textbf{\pa{ਫ਼ਰੇਮ}} (\emph{farem}) is borrowed from English \emph{frame}.
\end{cousinweb}

\begin{grammarlens}[title={What you've built}]
One more word for this chapter.
\end{grammarlens}

\subsection*{Your turn}
\begin{itemize}
\raggedright
  \item \emph{Say it:} \emph{farem}
  \item \emph{Say it:} \emph{farem}, once more
  \item \emph{Say it:} \emph{tasvīr laī ikk navāṁ farem kharīdo.} — Buy a new frame for the picture.
  \item \emph{Recall:} say the Punjabi for a photograph, then the Punjabi for a photo album, then say \emph{farem} again
  \item \emph{Write it:} copy \textbf{\pa{ਫ਼ਰੇਮ}} once, checking the whole word
\end{itemize}

\begin{tcolorbox}[breakable,colback=teal!4,colframe=teal!35!black,title={Before you move on}]
What does \pa{ਫ਼ਰੇਮ} mean? (\textquotedblleft{}A picture frame\textquotedblright{}.) Say it once more.
\end{tcolorbox}
